\chapter{Описание проекта}
\label{ch:project}

\section{Технологический стек}

Приложение построено \textbf{исключительно базовыми технологиями} браузера,
без сборщиков, транспиляторов и менеджеров зависимостей --- этого требует
методичка, и это же подчёркивает: для интерактивной страницы не нужен
<<ритуал>> современного фронтенда.

\begin{itemize}
  \item \textbf{HTML5} --- семантическая разметка резюме и панели управления;
  \item \textbf{CSS3} --- двухколоночная вёрстка на \texttt{CSS~Grid},
        анимации (\texttt{@keyframes}), адаптивность (\texttt{@media}) и
        отдельная печатная версия (\texttt{@media~print});
  \item \textbf{Vanilla JavaScript} (ES2024) --- рендер резюме, работа с
        \texttt{localStorage}, чтение фотографии, эффект Material Wave;
        классические теги \texttt{<script>} без модулей, поэтому страница
        работает и при открытии напрямую через \texttt{file://}.
\end{itemize}

Никаких \texttt{npm~install} и \texttt{node\_modules}: проект --- это набор
статических файлов, которые браузер исполняет как есть.

\section{Архитектура приложения}

Резюме не свёрстано вручную в \texttt{index.html}, а \textit{собирается из
данных}. Это разделяет содержимое и разметку и упрощает заполнение резюме.

\begin{lstlisting}[caption={Поток данных при загрузке страницы},captionpos=b]
config.js (RESUME_DATA)            index.html (пустой каркас)
        |                                   |
        +------> renderResume(data) ---------+
                        |
                        v
              applyFields(localStorage)  <-- правки пользователя
                        |
                        v
                 готовое резюме в DOM
\end{lstlisting}

Объект \texttt{RESUME\_DATA} из файла \texttt{config.js} описывает резюме:
имя, должность, контакты, навыки, языки, опыт, образование. Функция
\texttt{renderResume} строит из него разметку и вставляет в страницу. Затем
\texttt{applyFields} накладывает поверх сохранённые в \texttt{localStorage}
правки пользователя, если они есть.

Файлы JavaScript разделены по образцу учебного шаблона курса:
\begin{itemize}
  \item \texttt{config.js} --- только данные (объект \texttt{RESUME\_DATA});
  \item \texttt{resume.js} --- чистые функции: рендер, работа с хранилищем,
        чтение фотографии, создание <<волны>>, запуск анимаций;
  \item \texttt{start.js} --- оркестрация: функция \texttt{start()}
        навешивает обработчики событий по \texttt{DOMContentLoaded}.
\end{itemize}

\section{Конфиденциальность данных}

Репозиторий публичный, поэтому в файл \texttt{config.js} помещён
\textbf{выдуманный} персонаж --- фронтенд-разработчик с вымышленными ФИО,
контактами и нейтральным аватаром. Реальное резюме заполняется иначе:
пользователь открывает страницу и редактирует текст прямо в браузере, а
фотографию загружает кликом по аватару. Все реальные данные и фото хранятся
\textbf{только в \texttt{localStorage} браузера} и ни в один файл проекта не
записываются --- то есть физически не могут попасть в систему контроля
версий. Кнопка <<Сбросить>> очищает хранилище и возвращает страницу к
выдуманному шаблону.

\section{Структура каталога \texttt{lab14/}}

\begin{lstlisting}[caption={Файловая структура \texttt{lab14/}},captionpos=b]
lab14/
├── README.md
├── deploy.sh             # деплой/снос сайта за Caddy
├── resume/               # сайт-резюме
│   ├── index.html        # каркас + панель управления
│   ├── css/style.css     # вёрстка, анимации, печать, адаптив
│   ├── js/
│   │   ├── config.js     # RESUME_DATA — выдуманный персонаж
│   │   ├── resume.js     # чистые функции (рендер, storage, ripple)
│   │   └── start.js      # start() — навешивание событий
│   └── img/avatar.svg    # нейтральный аватар по умолчанию
└── latex-report/         # этот отчёт
\end{lstlisting}

\section{Способ публикации}

Готовое приложение публикуется на собственном домашнем сервере
(\texttt{Debian~12}, Docker Compose), на котором уже работает обратный прокси
\textbf{Caddy} с автоматическими TLS-сертификатами Let's~Encrypt. Сайт-резюме
добавляется в эту инфраструктуру как ещё один Docker-контейнер на образе
\texttt{nginx:alpine}, раздающий статические файлы каталога \texttt{resume/}.

\begin{lstlisting}[caption={Маршрут запроса до сайта},captionpos=b]
браузер  ->  443/tcp  ->  Caddy (Let's Encrypt TLS)
                              |
                              +-> resume:80 (nginx) -> /usr/share/nginx/html
\end{lstlisting}

Адрес сайта --- \url{https://resume.server34.netcraze.club}. Деплой
автоматизирован идемпотентным bash-скриптом \texttt{deploy.sh}; подробное
описание его шагов приведено в~\hyperref[ch:practice]{главе~3}. На сервере
оказывается резюме с выдуманным персонажем --- личные данные пользователя
остаются только в его браузере.

\endinput
